\chapter{Theoretische Grundlagen und Problemanalyse}

\section{Funktionsweise von \ac{LLM}-\acp{API}}

Die Kommunikation mit \ac{LLM}-\acp{API} erfolgt typischerweise über \ac{REST}ful \acp{API}, bei denen Anfragen
an zentrale Server der \ac{API}-Provider gesendet werden.
Die Anfragen enthalten dabei den sogenannten Prompt sowie verschiedene Parameter wie Temperatur, Maximalzahl an Tokens, usw.
Die Verbindung zum \ac{API}-Provider wird typischerweise mit \ac{TLS} verschlüsselt und verhindert somit
Man-in-the-Middle-Angriffe und das Abhören von Dritten.
Da jedoch unverschlüsselte Daten für die Verarbeitung der Modelle notwendig sind,
kann der Zugriff des \ac{API}-Providers auf die Informationen nicht verhindert werden (European Union Agency for Cybersecurity, 2022).
Trotz \acp{SLA}s, welche die Sicherheit der Daten garantieren sollen,
entsteht ein Vertrauensproblem zwischen \ac{API}-Provider und \ac{API}-Konsumierer,
da die Kontrolle über die Daten bei externen Parteien liegt (Kroll \& Zittrain, 2023). Dies stellt vor allem ein Problem dar,
wenn die \ac{API}-Provider unter einer anderen regulatorischen oder wettbewerbswirksamen Umgebung als der \ac{API}-Konsumierer stehen.


\section{Spezifische Risiken in der \ac{IT}-Infrastruktur}
Um die Bedeutung der Daten zu verstehen, lohnt es sich diese zu kategorisieren.
Es wird unterschieden zwischen personenbezogenen Daten (\ac{PII}) einerseits und
geschäftskritischen Infrastrukturdaten andererseits. Personenbezogene Daten umfassen dabei alle Daten, die eine natürliche Person eindeutig identifizieren können.
Das können zum Beispiel Namen, E-Mail-Adressen, Telefonnummern, etc. sein. Bei diesen greift die \ac{DSGVO},
weswegen eine unachtsame Verwendung der Daten zu erheblichen Geldbußen und Reputationsschäden führen kann.
Die Kategorie der geschäftskritischen Infrastrukturdaten umfasst hingegen alle Daten, die für die Betriebsführung relevant sind.
Dazu zählen in diesem Fall unter anderem IP-Adressen und Netzwerksegmentierungen, Servernamen und deren Betriebssystemversionen, etc.
Im Gegensatz zu personenbezogenen Daten existieren dabei keine regulatorischen Rahmenbedingungen, die die Verwendung der Daten einschränken.
Das Risiko liegt stattdessen darin, dass die Offenlegung dieser die Sicherheitsinfrastruktur des Unternehmens massiv
kompromittieren kann (Microsoft Research, 2023).